\section{Conclusion}
\label{sec:conclusion}

Behavioral deception detection in LLMs primarily detects instruction-following artifacts. Three controls---prompt equalization, cross-family extraction, regex baseline---collapse accuracies from 93.9--100\% to 52--69\%, with instruction-following contributing $+$7.5--15\,pp (14B/70B).

The 80.1\% label-free rule serves three purposes: (1) \textbf{calibration-free baseline}; (2) \textbf{extractor-independent ceiling}; (3) threshold for non-surface-lexical signal. Detectors <80.1\% under equalization lack validated behavioral signal. The construct-valid pipeline achieves 54.5\%; full cross-family 64.7\%. The signal does not transfer to fully-autonomous deception (four of six $n=200$ cells include chance); sycophancy (68.5/83\% at 3B/14B, $n\!=\!200$; 72\% at 70B, $n\!=\!50$ preliminary) is the one positive case. Qwen~32B collapse ($n\!=\!100$, $\mu_\mathrm{lie}\!=\!\mu_\mathrm{truth}\!=\!0.00$) is Qwen-family-specific: cross-organizational replication on Gemma~2~27B ($n\!=\!100$, 84\%, $p\!<\!0.001$) and Mistral~7B ($n\!=\!100$, 56\%, ns) confirms the zero-marker collapse does not generalize to other heavily-aligned families.

\paragraph{Reproducibility.} Code and data are released as supplementary material. All experiments use publicly available models via AWS Bedrock or HuggingFace; open-weight results require no API keys.

All refusal-marker patterns are English-only. Future work should apply the three-control protocol to other behavioral detectors; knowledge-probing approaches that elicit correction behavior directly may prove more robust to RLHF smoothing.
